\section{Discussion, threats and related work}
\label{sec:discussion}

\textbf{What this means for agent builders and reviewers.} A performance
PR whose mechanism leaves no kernel footprint, or whose speed-up comes with
a changed output, should not be merged on its description. Kernel signals
are cheap to collect once a workload exists; asking agents to attach the
trace evidence for the mechanism they claim would turn many unverifiable
claims into checkable ones.

\textbf{Threats to validity.} The measured sample is small and early
(\nAnalysed PRs, up to \maxPerRepo from one repository), so the shares in
Table~\ref{tab:quadrant} are preliminary. Workloads are written by an
agent: although it never sees the claim and an independent check enforces
coverage of the changed lines, a workload may not reflect real use. Because
the agent is blind to the claim, it may exercise a different part of a
large PR than the one the claim is about; the review removed such PRs when
it found them, but others may remain and would bias towards
``mechanism absent''.
Classification and claim extraction are done by LLMs; a human check of a
stratified sample of 300 PRs is prepared but not yet done. Kernel signals are evidence for a
mechanism, not proof; the control programs show that each signal responds
as expected. Each PR is measured once on one machine; repeated measurements on a
second machine may be added later.

\textbf{Related work.} Studies of agent performance PRs
characterize topics~\cite{opu2025bertopic}, optimization
patterns~\cite{peng2025optimization}, mutation
patterns~\cite{dovhoshliubnyi2026mutation} and how agents report
validation~\cite{peng2026validation}, all from text and code. Qi et
al.\ re-execute 53 fixes from AIDev~\cite{qi2026merged}. Gong et al.\
find agent PR descriptions that claim changes the code does not
implement~\cite{gong2026inconsistency}; we check a claim's runtime
effect instead of its presence in the diff. Benchmarks
such as PIE~\cite{shypula2024pie}, EffiBench~\cite{huang2024effibench},
SWE-Perf~\cite{he2025sweperf}, GSO~\cite{shetty2025gso},
PerfBench~\cite{garg2025perfbench}, SWE-fficiency~\cite{ma2025swefficiency}
and FormulaCode~\cite{sehgal2026formulacode} measure whether models or
agents can make given code faster, and their reliability has been
questioned~\cite{chen2026benchmarks}. Performance bugs in human-written
code have long been studied~\cite{jin2012perfbugs,zaman2012perfbugs}, and
kernel traces have been used to explain where execution time
goes~\cite{giraldeau2016wait,ezzatijivan2017multiscale}. We instead check,
with kernel traces, the mechanism that agent PRs claim.

\textbf{Conclusion.} Agents' performance claims are mostly unmeasured, and
a stopwatch cannot tell whether they are true for the stated reason. Our
pipeline checks the claimed mechanism in the kernel; the full study is
under way on \nJobs PRs.
